% GENERATED FILE. Edit canonical lessons, then run npm run generate:books.
% canonical-source-hash: fnv1a64:8c881db5ec83ca6e
% canonical-lessons: JA-C426-kibun, JA-C426-fuan, JA-C426-heiki, JA-C426-sukkiri, JA-C426-taeru

\chapter{How I Feel Today}
\label{ch:ja-how-i-feel-today}
\hlchaptermodality{426}

\begin{chapteropening}
\textbf{By the end of this chapter:} \emph{I can say how I feel and what I put up with.}

Complete the last lesson of chapter 426: I can say how I feel and what I put up with.
\end{chapteropening}

\section[\texorpdfstring{\ja{きぶん} (kibun)}{きぶん (kibun)}]{\ja{きぶん} (kibun) — one's mood; how one feels}

\label{lesson:JA-C426-kibun}



\begin{quote}
Before the new one: say the Japanese for serious, then the Japanese for selfish.
\end{quote}

\subsection*{You'll want to know: \ja{きぶん}}
\textbf{\ja{きぶん}} — \emph{kibun} — \textquotedblleft{}one's mood; how one feels\textquotedblright{}.

\begin{quote}
\textbf{\ja{けさは} \ja{きぶんが} \ja{いいです。}} — \emph{kesa wa kibun ga ii desu} — \textquotedblleft{}I am in a good mood this morning.\textquotedblright{}
\end{quote}

\begin{cousinweb}
\textbf{\ja{きぶん}} (\emph{kibun}) is Sino-Japanese (kango): two Chinese characters read with their on'yomi, readings borrowed from Chinese. They are \emph{ki}, \textquotedblleft{}spirit\textquotedblright{} (Middle Chinese \emph{khj+jH}) and \emph{bun}, \textquotedblleft{}part\textquotedblright{} (\emph{pjun}) (KANJIDIC2; Middle Chinese in Baxter's transcription).
\end{cousinweb}

\begin{grammarlens}[title={What you've built}]
One more word for describing what you lived and felt.
\end{grammarlens}

\subsection*{Your turn}
\begin{itemize}
\raggedright
  \item \emph{Say it:} \emph{kibun}
  \item \emph{Say it:} \emph{kibun}, once more
  \item \emph{Say it:} \emph{kesa wa kibun ga ii desu} — I am in a good mood this morning.
  \item \emph{Recall:} say the Japanese for serious, then the Japanese for selfish, then say \emph{kibun} again
\end{itemize}

\begin{tcolorbox}[breakable,colback=teal!4,colframe=teal!35!black,title={Before you move on}]
What does \ja{きぶん} mean? (\textquotedblleft{}One's mood; how one feels\textquotedblright{}.) Say it once more.
\end{tcolorbox}
\section[\texorpdfstring{\ja{ふあん} (fuan)}{ふあん (fuan)}]{\ja{ふあん} (fuan) — anxious; uneasy}

\label{lesson:JA-C426-fuan}



\begin{quote}
Before the new one: say the Japanese for selfish, then the Japanese for one's mood.
\end{quote}

\subsection*{You'll want to know: \ja{ふあん}}
\textbf{\ja{ふあん}} — \emph{fuan} — \textquotedblleft{}anxious; uneasy\textquotedblright{}.

\begin{quote}
\textbf{\ja{しけんの} \ja{まえは} \ja{ふあんでした。}} — \emph{shiken no mae wa fuan deshita} — \textquotedblleft{}I felt uneasy before the exam.\textquotedblright{}
\end{quote}

A na-adjective: \ja{な} before a noun.

\begin{cousinweb}
\textbf{\ja{ふあん}} (\emph{fuan}) is Sino-Japanese (kango): two Chinese characters read with their on'yomi, readings borrowed from Chinese. They are \emph{fu}, \textquotedblleft{}negative\textquotedblright{} (Middle Chinese \emph{pjuw}) and \emph{an}, \textquotedblleft{}peaceful\textquotedblright{} (\emph{'an}) (KANJIDIC2; Middle Chinese in Baxter's transcription).
\end{cousinweb}

\begin{grammarlens}[title={What you've built}]
One more word for describing what you lived and felt.
\end{grammarlens}

\subsection*{Your turn}
\begin{itemize}
\raggedright
  \item \emph{Say it:} \emph{fuan}
  \item \emph{Say it:} \emph{fuan}, once more
  \item \emph{Say it:} \emph{shiken no mae wa fuan deshita} — I felt uneasy before the exam.
  \item \emph{Recall:} say the Japanese for selfish, then the Japanese for one's mood, then say \emph{fuan} again
\end{itemize}

\begin{tcolorbox}[breakable,colback=teal!4,colframe=teal!35!black,title={Before you move on}]
What does \ja{ふあん} mean? (\textquotedblleft{}Anxious; uneasy\textquotedblright{}.) Say it once more.
\end{tcolorbox}
\section[\texorpdfstring{\ja{へいき} (heiki)}{へいき (heiki)}]{\ja{へいき} (heiki) — unbothered; all right}

\label{lesson:JA-C426-heiki}



\begin{quote}
Before the new one: say the Japanese for one's mood, then the Japanese for anxious.
\end{quote}

\subsection*{You'll want to know: \ja{へいき}}
\textbf{\ja{へいき}} — \emph{heiki} — \textquotedblleft{}unbothered; all right\textquotedblright{}.

\begin{quote}
\textbf{\ja{ちちは} \ja{さむい} \ja{あさも} \ja{へいきです。}} — \emph{chichi wa samui asa mo heiki desu} — \textquotedblleft{}My father does not mind even cold mornings.\textquotedblright{}
\end{quote}

A na-adjective: \ja{な} before a noun.

\begin{cousinweb}
\textbf{\ja{へいき}} (\emph{heiki}) is Sino-Japanese (kango): two Chinese characters read with their on'yomi, readings borrowed from Chinese. They are \emph{hei}, \textquotedblleft{}peace\textquotedblright{} (Middle Chinese \emph{bjaeng}) and \emph{ki}, \textquotedblleft{}spirit\textquotedblright{} (\emph{khj+jH}) (KANJIDIC2; Middle Chinese in Baxter's transcription).
\end{cousinweb}

\begin{grammarlens}[title={What you've built}]
One more word for describing what you lived and felt.
\end{grammarlens}

\subsection*{Your turn}
\begin{itemize}
\raggedright
  \item \emph{Say it:} \emph{heiki}
  \item \emph{Say it:} \emph{heiki}, once more
  \item \emph{Say it:} \emph{chichi wa samui asa mo heiki desu} — My father does not mind even cold mornings.
  \item \emph{Recall:} say the Japanese for one's mood, then the Japanese for anxious, then say \emph{heiki} again
\end{itemize}

\begin{tcolorbox}[breakable,colback=teal!4,colframe=teal!35!black,title={Before you move on}]
What does \ja{へいき} mean? (\textquotedblleft{}Unbothered; all right\textquotedblright{}.) Say it once more.
\end{tcolorbox}
\section[\texorpdfstring{\ja{すっきり} (sukkiri)}{すっきり (sukkiri)}]{\ja{すっきり} (sukkiri) — refreshed; clear and tidy}

\label{lesson:JA-C426-sukkiri}



\begin{quote}
Before the new one: say the Japanese for anxious, then the Japanese for unbothered.
\end{quote}

\subsection*{You'll want to know: \ja{すっきり}}
\textbf{\ja{すっきり}} — \emph{sukkiri} — \textquotedblleft{}refreshed; clear and tidy\textquotedblright{}.

\begin{quote}
\textbf{\ja{よく} \ja{ねたので、きぶんが} \ja{すっきりしました。}} — \emph{yoku neta node, kibun ga sukkiri shimashita} — \textquotedblleft{}I slept well, so my head feels clear.\textquotedblright{}
\end{quote}

\begin{cousinweb}
\textbf{\ja{すっきり}} (\emph{sukkiri}) is a native Japanese mimetic word, inherited rather than borrowed from Chinese, and its older origin is not settled (JMdict).
\end{cousinweb}

\begin{grammarlens}[title={What you've built}]
One more word for describing what you lived and felt.
\end{grammarlens}

\subsection*{Your turn}
\begin{itemize}
\raggedright
  \item \emph{Say it:} \emph{sukkiri}
  \item \emph{Say it:} \emph{sukkiri}, once more
  \item \emph{Say it:} \emph{yoku neta node, kibun ga sukkiri shimashita} — I slept well, so my head feels clear.
  \item \emph{Recall:} say the Japanese for anxious, then the Japanese for unbothered, then say \emph{sukkiri} again
\end{itemize}

\begin{tcolorbox}[breakable,colback=teal!4,colframe=teal!35!black,title={Before you move on}]
What does \ja{すっきり} mean? (\textquotedblleft{}Refreshed; clear and tidy\textquotedblright{}.) Say it once more.
\end{tcolorbox}
\section[\texorpdfstring{\ja{たえる} (taeru)}{たえる (taeru)}]{\ja{たえる} (taeru) — to endure; to bear}

\label{lesson:JA-C426-taeru}



\begin{quote}
Before the new one: say the Japanese for unbothered, then the Japanese for refreshed.
\end{quote}

\subsection*{You'll want to know: \ja{たえる}}
\textbf{\ja{たえる}} — \emph{taeru} — \textquotedblleft{}to endure; to bear\textquotedblright{}.

\begin{quote}
\textbf{\ja{あには} \ja{あつい} \ja{なつに} \ja{たえています。}} — \emph{ani wa atsui natsu ni taete imasu} — \textquotedblleft{}My elder brother is putting up with the hot summer.\textquotedblright{}
\end{quote}

Polite \textbf{\ja{たえます}}, \emph{taemasu}; an ichidan verb. Past (polite): \textbf{\ja{ははは} \ja{いたみに} \ja{たえました。}}, \emph{haha wa itami ni taemashita}.

\begin{cousinweb}
\textbf{\ja{たえる}} (\emph{taeru}) is a native Japanese word (wago), inherited rather than borrowed from Chinese; its kanji spelling was attached to it later (JMdict).
\end{cousinweb}

\begin{grammarlens}[title={What you've built}]
One more word for describing what you lived and felt.
\end{grammarlens}

\subsection*{Your turn}
\begin{itemize}
\raggedright
  \item \emph{Say it:} \emph{taeru}
  \item \emph{Say it:} \emph{taeru}, once more
  \item \emph{Say it:} \emph{ani wa atsui natsu ni taete imasu} — My elder brother is putting up with the hot summer.
  \item \emph{Recall:} say the Japanese for unbothered, then the Japanese for refreshed, then say \emph{taeru} again
\end{itemize}

\begin{tcolorbox}[breakable,colback=teal!4,colframe=teal!35!black,title={Before you move on}]
What does \ja{たえる} mean? (\textquotedblleft{}To endure; to bear\textquotedblright{}.) Say it once more.
\end{tcolorbox}
